% American Mathematical Society article.
\documentclass[11pt]{amsart}

\usepackage{amssymb}

\newtheorem{theorem}{Theorem}[section]
\newtheorem{lemma}[theorem]{Lemma}
\theoremstyle{definition}
\newtheorem{definition}[theorem]{Definition}
\theoremstyle{remark}
\newtheorem{remark}[theorem]{Remark}

\begin{document}

\title{An Example Title}

\author{A. U. Thor}
\address{Department of Mathematics, Your University, City, Country}
\email{author@example.org}

\subjclass[2020]{Primary 00A00; Secondary 00B00}
\keywords{First keyword, second keyword}

\begin{abstract}
State the main result and the idea of its proof in a few sentences, with no
citations.
\end{abstract}

\maketitle

\section{Introduction}

Say what is known, what is missing, and what this paper proves. Cite earlier
work by key \cite{einstein1905,knuth1984}.

\section{Main result}

\begin{definition}
A sequence $(a_n)$ of real numbers is \emph{bounded} if there is an $M$ with
$|a_n| \le M$ for all $n$.
\end{definition}

\begin{theorem}\label{thm:main}
Every bounded monotone sequence of real numbers converges.
\end{theorem}

\begin{proof}
Let $(a_n)$ be increasing and bounded, and let $L = \sup_n a_n$. Given
$\varepsilon > 0$, some $a_N$ exceeds $L - \varepsilon$, and then
$L - \varepsilon < a_n \le L$ for every $n \ge N$.
\end{proof}

\begin{remark}
Theorem~\ref{thm:main} fails over the rational numbers.
\end{remark}

\bibliographystyle{amsplain}
\bibliography{references}

\end{document}
